%%    Process

\subsection{Glyph: \glyph{Process}}
\label{sec:af:process}

In \SBGNAFLone, a \emph{process} is used to represent a conceptual outcome of a sequential series of actions, motions, or occurrences, such as chemical reactions, that usually involve multiple biological activities from multiple entities.

Although a \emph{process} may also have implications of biological activities, it is different from the glyph \emph{biological activity} in that the latter is the activity from a defined EPN.



\begin{glyphDescription}

\glyphSboTerm SBO:0000375 ! process

\glyphContainer A \glyph{process} is represented by a double-lined rectangle, as shown in \fig{af:process}. 

\glyphLabel A \glyph{process} is identified by a label placed in an unbordered box containing a string of characters.  The characters can be distributed on several lines to improve readability, although this is not mandatory.  The label box must be attached to the center of the container.  The label may spill outside of the container.

\glyphAux None.

\end{glyphDescription}

\begin{figure}[H]
  \centering
  \includegraphics[scale=0.1]{images/process}
  \caption{The \AF glyph for \glyph{process}.}
  \label{fig:af:process}
\end{figure}
